\section{Several minimizers}\label{sec:optsets}
The growth analysis of Section~\ref{sec:growth} localizes every retained
coordinate near one minimizer. When the optimal set $\mathcal S$ is not a
single point, this localization fails, even though Theorem~\ref{thm:transfer}
shows that exactness itself costs only $\poly(I)+\log_2(1/g_S)$ bits of
accuracy under growth $g_S$ towards the set. This section treats three
situations: finitely many minimizers, where unions of uniform cells restore a
bound that is polynomial for fixed $p$; coordinatewise concave quadratics,
where the decomposition describes the entire optimal set; and the class of
instances certified by a diagonal Lagrangian multiplier, where an unknown
optimal set is found exactly. Throughout, $F=\frac12x^\T Hx+b^\T x+c$ is a
rational quadratic, and a number $g_S>0$ is a \emph{set-growth constant} if
\begin{equation}\label{eq:setgrowth}
F(x)-\OPT\ge g_S\dist(x,\mathcal S)^2\qquad(x\in X);
\end{equation}
then $\kappa_S=\max\{1,L/g_S\}$. For a bounded mixed box some $g_S>0$ always
exists (Remark~\ref{rem:setgrowth}), but its size is not controlled, and
nothing below computes it.

\subsection{One center cannot serve several minimizers}\label{sec:exact-nonunique}

Theorem~\ref{thm:transfer} reduces exact output to reaching accuracy
$\varepsilon_S$. The algorithm of Section~\ref{sec:growth} does not reach it
at a rate controlled by $(p,\kappa)$ when the minimizer is not unique, even
with two isolated minimizers and set growth: its grids are graded around one
center, and coarse cells near a distant minimizer keep the corrected minimum
low.

\begin{proposition}[One center cannot certify two distant minimizers]
\label{prop:twocenters}
For $M>0$ let
\[
F_M(x,z)=x^2-2xz+Mz\qquad\text{on }X=[0,M]^2\ (\text{both coordinates
continuous}).
\]
Then $S=\{(0,0),(M,M)\}$, $\OPT=0$, $L_x=2$, $L_z=0$, and
$F_M(v)\ge\frac1{20}\dist(v,S)^2$ on $X$, so the set-growth condition number
is at most $40$; one bag $\{x,z\}$ covers both scopes ($p=2$).
Let $\theta\in(0,1/4]$, $h\in(0,M]$ and let $c\in X$ be any center. For the
graded grids~\eqref{eq:step} on $X$ with center $c$, mesh $h$, grading
$\theta$, and any corrections $d_x=L_xw_x^2/8$ with $L_x\ge2$ and $d_z\ge0$,
\begin{equation}\label{eq:twocenters}
\beta=\min_GQ\le-\max\Bigl\{\frac{\theta^2M^2}{379},\ \frac{h^2}{20}\Bigr\}.
\end{equation}
Consequently, in FG and CT every stage has certified gap
$U-\beta\ge\max\{\theta^2M^2/379,\,h^2/20\}$; a stage certifying a
gap $\varepsilon$ uses grids with at least $M/(48\sqrt\varepsilon)$ nodes in
each coordinate; and any exact-output procedure based on these certificates
that accepts only when $U-\beta<1$, such as EX, builds a table with at least
$M^2/2304$ entries. For integer $M$ the input length is $I=O(\log M)$, so
this work is $2^{\Omega(I)}$.
\end{proposition}

\begin{proof}
\emph{The instance.} $F_M=x^2+z(M-2x)$. For $x<M/2$ the minimum over $z$ is
at $z=0$ with value $x^2$; for $x>M/2$ it is at $z=M$ with value $(M-x)^2$;
at $x=M/2$ the value is $M^2/4$. So $S$ and $\OPT$ are as stated, and the
Hessian $\bigl(\begin{smallmatrix}2&-2\\-2&0\end{smallmatrix}\bigr)$ gives
$L_x=2$, $L_z=0$. The map $(x,z)\mapsto(M-x,M-z)$ preserves $F_M$ (a direct
expansion) and swaps the minimizers, so for growth we may assume
$x\le M/2$. If $x\le M/4$, then $M-2x\ge M/2\ge z/2$ and
$F_M\ge x^2+z^2/2\ge\frac12\norm{(x,z)}^2$. If $M/4<x\le M/2$, then
$F_M\ge x^2\ge M^2/16$ while $\norm{(x,z)}^2\le5M^2/4$. In both cases
$F_M\ge\frac1{20}\dist((x,z),S)^2$.

\emph{The bound~\eqref{eq:twocenters}.} By symmetry assume $M-c_x\ge M/2$
(otherwise use the reflected argument with $z=0$ and the endpoint $0$). Let
$c_x=v_0<v_1<\dots<v_K=M$ be the $x$-nodes on the right of the center,
$D_k=v_k-c_x$, $a=v_{K-1}-v_{K-2}$ (if $K\ge2$) and $c'=M-v_{K-1}$. All steps
except the last are unclipped: $D_{k+1}=(1+\theta)D_k+h$ for $k<K-1$, and
$D_K\le(1+\theta)D_{K-1}+h$. The node $z=M$ is in the $z$-grid, and
$F_M(v,M)=(M-v)^2$. Every correction is nonnegative and
$d_x(v)\ge w_x(v)^2/4$.

If $K=1$, the node $M$ has the adjacent interval $[c_x,M]$ of length at least
$M/2$, so $\beta\le Q(M,M)\le-M^2/16$, and
$M^2/16\ge\max\{\theta^2M^2/379,h^2/20\}$ since $\theta\le1/4$ and $h\le M$.

If $K\ge2$, then $Q(M,M)\le-c'^2/4$ and
$Q(v_{K-1},M)\le c'^2-a^2/4$. If $c'^2\ge a^2/5$ the first is at most
$-a^2/20$, otherwise the second is below $-a^2/20$; so $\beta\le-a^2/20$.
Now $a=h+\theta D_{K-2}\ge h$, and the recursion gives
$D_{K-2}\ge(D_K-(2+\theta)h)/(1+\theta)^2$. If $h\le M/16$, then with
$D_K\ge M/2$ and $\theta\le1/4$,
$D_{K-2}\ge(\frac M2-\frac{9M}{64})\cdot\frac{16}{25}=\frac{23M}{100}$, so
$a\ge\frac{23}{100}\theta M$. If $h>M/16$, then $a\ge h>M/16\ge\theta M/4$.
Hence $a\ge\max\{h,\frac{23}{100}\theta M\}$ and
$\beta\le-\max\{h^2,(0.23\,\theta M)^2\}/20$, which
implies~\eqref{eq:twocenters} because $0.23^2/20>1/379$.

\emph{Consequences.} Filtering keeps both minimizers
(Proposition~\ref{prop:filter}), so every box of FG is $[0,M]^2$, and every
stage is of the form above with $U\ge\OPT=0$; this gives the gap bound. A
stage with $U-\beta\le\varepsilon$ has $\theta\le\sqrt{379\varepsilon}/M$
and $h\le\sqrt{20\varepsilon}$, so every step is at most
$h+\theta M<24\sqrt\varepsilon$, and the side of length at least $M/2$ in each
coordinate needs at least $M/(48\sqrt\varepsilon)$ nodes. The table of the
bag $\{x,z\}$ has at least $M^2/(2304\varepsilon)$ entries. CT stops only
after some stage certifies its target, and an acceptance with
$U-\beta<1$ needs a stage with gap below one.
\end{proof}

The instance is trivial for other methods; the point is that the failure
comes from the single center, not from conditioning. Replacing hulls by
unions of retained cells, on uniform meshes, removes it at the price of a
$\sqrt n$ factor per coordinate.

\paragraph{Algorithm UC$(\varepsilon)$ (uniform cells).}
For each coordinate keep a finite set $\mathcal C_i$ of \emph{retained cells},
closed intervals of positive length; initially $\mathcal C_i=\{[\ell_i,u_i]\}$.
Let $J$ be the least $j\ge0$ with $\frac12nLs^24^{-j}\le\varepsilon$, start
with $\hat x=\ell$, $U=F(\ell)$, and for $j=0,1,\dots,J$ with $h_j=s2^{-j}$:
\begin{enumerate}[label=(\roman*)]
\item \emph{Refine.} Let $\Lambda_{ij}=\{\ell_i+kh_j:k\in\Z_{\ge0}\}$ for
$i\in I_C$ and $\Lambda_{ij}=\{\ell_i+\lceil kh_j\rceil:k\in\Z_{\ge0}\}$ for
$i\in I_Z$. A retained cell $[a,b']$ is split at the points of
$\Lambda_{ij}\cap(a,b')$; the resulting \emph{stage cells} have consecutive
nodes as endpoints, and $G_i$ is the set of all their endpoints. Let
$w_i(v)$ be the largest length of a stage cell with endpoint $v$, ignoring
integer cells of length one, and define $d_i$, $\Delta$, $Q$, $\beta_j$ and $m_i$
as in~\eqref{eq:correction}--\eqref{eq:minmarginal} on $G=\prod_iG_i$.
\item \emph{Solve.} Compute $\beta_j$, a minimizer $y_j$ and all $m_i$ by
Lemma~\ref{lem:dp}; if $F(y_j)<U$, set $\hat x=y_j$, $U=F(y_j)$.
\item \emph{Filter.} Let $\mathcal C_i$ be the set of stage cells $[a,b']$
with $\min\{m_i(a),m_i(b')\}\le U$. Do not take hulls.
\end{enumerate}
Return $\hat x$, $U$ and $\beta_J$. The current domain at stage $j$ is
$X_j=X\cap\prod_i\bigcup\mathcal C_i$.

Since $\Lambda_{i,j-1}\subseteq\Lambda_{ij}$ (for integers,
$\lceil kh_{j-1}\rceil=\lceil 2kh_j\rceil$), every stage cell is either a
segment between consecutive points of $\Lambda_{ij}\cap[\ell_i,u_i]$ or such
a segment clipped at $u_i$. Hence a continuous stage cell has length at most
$h_j$; an integer stage cell has length at most $\lceil h_j\rceil$, and if
that length is at least two then $h_j>1$ and $\lceil h_j\rceil<2h_j$.
Consecutive points of $\Lambda_{ij}$ are at least $h_j$ apart for $i\in I_C$
and at least $\max\{1,\lfloor h_j\rfloor\}\ge h_j/2$ apart for $i\in I_Z$.
Each stage cell contains at most one point of $\Lambda_{i,j+1}$ in its
interior.

\begin{lemma}[Cells: interpolation and safe filtering]\label{lem:cells}
Assume only Definition~\ref{def:curvature}. At every stage, for every
$x\in X_j$ there is a random $Y\in G$ with $\E Q(Y)\le F(x)$ and
$Y_i\in\{a,b'\}$ whenever $x_i$ lies in the stage cell $[a,b']$. Hence
$\beta_j\le\min_{X_j}F$, and $\min\{m_i(a),m_i(b')\}\le F(x)$ for every
$x\in X_j$ with $x_i\in[a,b']$. Every global minimizer lies in every $X_j$,
and $\beta_J\le\OPT\le U$.
\end{lemma}

\begin{proof}
For each $i$ choose a stage cell $[a_i,b_i]$ containing $x_i$. If $x_i$ is
an endpoint put $Y_i=x_i$; otherwise let $Y_i=b_i$ with probability
$(x_i-a_i)/(b_i-a_i)$ and $Y_i=a_i$ otherwise, independently. If $i\in I_Z$
and $b_i-a_i=1$, then $x_i$ is an endpoint. Let
$Z^{k}=(Y_1,\dots,Y_k,x_{k+1},\dots,x_n)\in\bar X$. Conditionally on
$Y_1,\dots,Y_{k-1}$, the function $\psi(t)=F(Z^{k-1}+(t-x_k)e_k)-\frac{L_k}2t^2$
is concave on $[a_k,b_k]$, so Jensen's inequality gives
$\E F(Z^k)\le F(Z^{k-1})+\frac{L_k}2\Var(Y_k)$. Summing over $k$,
$\E F(Y)\le F(x)+\sum_k\frac{L_k}2\Var(Y_k)$. If $Y_k$ is rounded, then
$\Var(Y_k)=(x_k-a_k)(b_k-x_k)\le(b_k-a_k)^2/4$ while both possible values have
$d_k\ge L_k(b_k-a_k)^2/8$; if not, $\Var(Y_k)=0\le d_k$. So
$\E D(Y)\ge\sum_k\frac{L_k}2\Var(Y_k)$ and $\E Q(Y)\le F(x)$. Since
$Q(Y)\ge\beta_j$ and, when $x_i\in[a,b']$ is the chosen cell, $Q(Y)\ge
\min\{m_i(a),m_i(b')\}$, the two bounds follow. A point $x\in X_j$ removed
at stage $j$ has some coordinate whose stage cells containing $x_i$ were all
discarded, so $F(x)>U_j\ge U$; in particular minimizers are never removed,
and by induction $X_J$ contains $S$ and $\beta_J\le\OPT$.
\end{proof}

\begin{theorem}[Uniform cells]\label{thm:cells}
Assume Definition~\ref{def:curvature}. For every $\varepsilon>0$,
UC$(\varepsilon)$ returns a feasible $\hat x$ and $\beta\le\OPT\le
F(\hat x)\le\beta+\varepsilon$ with a certificate (all stage tables and
filtering decisions) whose validity uses only Definition~\ref{def:curvature}.
Suppose moreover that $S$ is finite, $r=\max_i\abs{\pi_i(S)}$, and
$F(x)-\OPT\ge g_S\dist(x,S)^2$ on $X$; put $\kappa_S=\max\{1,L/g_S\}$. Then
every grid of every stage has at most
\begin{equation}\label{eq:cellcount}
K_S=12\,r\,\bigl(2\sqrt{n\kappa_S}+1\bigr)
\end{equation}
nodes per coordinate. For an explicit rational quadratic, UC$(2^{-q})$ uses
$O\bigl(p(N+\abs{\mathcal A})K_S^p\bigr)\poly(I+q)$ bit operations, and
EX with UC in place of CT returns an exact minimizer in
$O\bigl(p(N+\abs{\mathcal A})K_S^p\bigr)\poly(I+\log\kappa_S)$ bit
operations. None of $g_S$, $S$, $r$ is an input.
\end{theorem}

\begin{proof}
\emph{Gap.} Every stage cell that receives a correction has length at most
$h_j$ (continuous) or below $2h_j$ (integer), so
$D(y)\le\sum_iL_i(2h_j)^2/8\le\frac12nLh_j^2$ for every $y\in G$. Hence
$U-\beta_j\le F(y_j)-\beta_j=D(y_j)\le\frac12nLh_j^2$, which is at most
$\varepsilon$ at $j=J$. Validity is Lemma~\ref{lem:cells}.

\emph{Count.} Fix a stage $j$ and a cell $[a,b']$ retained by step~(iii). One
endpoint $v$ has $m_i(v)\le U$; let $z\in G$ attain $m_i(v)$, so $z_i=v$.
By Lemma~\ref{lem:cells}, $\beta_j\le\OPT$, so
$U\le F(y_j)\le\OPT+\frac12nLh_j^2$ and
\[
F(z)-\OPT\le U+D(z)-\OPT\le nLh_j^2 .
\]
Set growth gives $\dist(z,S)^2\le nLh_j^2/g_S\le n\kappa_Sh_j^2$, so
$\abs{v-\sigma}\le\rho:=\sqrt{n\kappa_S}\,h_j$ for some $\sigma\in\pi_i(S)$.
The endpoints of stage cells are points of $\Lambda_{ij}$ or $u_i$. An
interval of length $2\rho$ contains at most $2\rho/\delta+1$ points of
$\Lambda_{ij}$, where $\delta=h_j$ (continuous) or $h_j/2$ (integer), plus
possibly $u_i$; each endpoint belongs to at most two stage cells. So at most
$2r(4\sqrt{n\kappa_S}+2)$ cells are retained, and since each contains at most
one new interior node, the stage-$(j+1)$ grid has at most
$3\cdot2r(4\sqrt{n\kappa_S}+2)=K_S$ nodes. Stage $0$ has two nodes per
coordinate.

\emph{Work and bit length.} By Lemma~\ref{lem:dp} each stage takes
$O(p(N+\abs{\mathcal A})K_S^p)$ table operations, and
$J=O(\log(nLs^2/\varepsilon))=O(I+q)$. Continuous nodes have the form
$\ell_i+ks2^{-j}$ and lie in $[\ell_i,u_i]$, integer nodes are integers, so
all nodes have denominators dividing $\Delta 2^j$ and bit length $O(I+j)$; factor
values, corrections and messages have bit length polynomial in $I+j$. For
exact output, Theorem~\ref{thm:transfer}(a) is reached at
$q\le\poly(I)+\log_2\kappa_S$, as in the proof of Theorem~\ref{thm:exact}.
\end{proof}

\begin{remark}[What grading buys]\label{rem:grading}
With $r=1$, Theorem~\ref{thm:cells} gives accuracy-independent grids of
size $O(\sqrt{n\kappa})$ without trial gradings or caps. The graded
single-center grids of Section~\ref{sec:growth} reduce $\sqrt n$ to
$\log n$, which is what makes Theorem~\ref{thm:approx} fixed-parameter
tractable in $(p,\kappa)$; Proposition~\ref{prop:twocenters} shows that
they pay for this with a dependence on a unique minimizer. Whether exact
output admits an $f(p,\kappa_S,r)\poly(I)$ bound when $S$ is finite, or any
accuracy-independent grid bound when $S$ is a continuum, is open.
\end{remark}

\begin{remark}[Relation to rational-witness results]\label{rem:np}
Corollary~\ref{cor:height} is the box case, with an explicit constant, of a
classical fact: a quadratic program with rational data that attains its
minimum has a global minimizer of polynomial encoding length, obtained from
a rational linear system \cite{Vavasis1990}; Vavasis used it to show that the
decision version of quadratic programming is in NP, and Del Pia, Dey and
Molinaro extended NP membership to mixed-integer quadratic programs with
unbounded integer variables \cite{DelPiaDeyMolinaro2017}. On a bounded
mixed box the integer coordinates have polynomial length, so that extension
is not needed. These results certify $\OPT\le t$. Deciding $\OPT\le t$ for
box QP is NP-complete, and NP-hard already at treewidth two
\cite{DelPiaKhajavirad2026}, so polynomial-size, polynomial-time checkable
certificates of lower bounds $\OPT\ge t$ for all instances would place an
NP-complete problem in coNP. The
certificates of Theorem~\ref{thm:exact} have size $f(p,\kappa)\poly(I)$:
they exist for bounded $(p,\kappa)$ and are checked without trusting
growth. The explicit constants $R$ and $\Omega$ are what make the stopping test
computable.
\end{remark}


\subsection{Coordinatewise concave quadratics: every optimizer at once}\label{sec:endpointset}

Write $E_i=\{l_i,u_i\}$ and $E(B)=\prod_{i\in B}E_i$. For $x\in\bar X$ and
$v_i\in E_i$ put
\[
e_i(v_i;x_i)=\begin{cases}u_i-x_i,&v_i=l_i,\\ x_i-l_i,&v_i=u_i,\end{cases}
\qquad
\pi_t(v;x)=\prod_{i\in B_t}\frac{e_i(v_i;x_i)}{u_i-l_i}\quad(v\in E(B_t)).
\]
Thus $\pi_t(\cdot;x)$ is the law of $Y_{B_t}$ when every coordinate of $x$
is rounded independently to an endpoint with mean $x_i$. Root $T$ at $t_0$,
let $\mathrm{ch}(t)$ be the children of $t$, and write
$B_t^{\uparrow}=B_t\cap B_{\mathrm{pa}(t)}$, with $B_{t_0}^{\uparrow}=\emptyset$.
Let $\phi_t$ be the sum of the monomials assigned to $t$.

\begin{lemma}[Endpoint Bellman residual identity]\label{lem:endpointid}
Assume $H_{ii}\le0$ for every $i$. Compute, in post-order and with all
arguments ranging over endpoint labels,
\begin{align*}
M_t(v_{B_t^\uparrow})&=\min_{v_{B_t\setminus B_t^\uparrow}}
\Bigl[\phi_t(v_{B_t})+\sum_{t'\in\mathrm{ch}(t)}M_{t'}(v_{B_{t'}^\uparrow})\Bigr],\\
r_t(v_{B_t})&=\phi_t(v_{B_t})+\sum_{t'\in\mathrm{ch}(t)}M_{t'}(v_{B_{t'}^\uparrow})
-M_t(v_{B_t^\uparrow})\ \ge0 .
\end{align*}
Then the scalar $M_{t_0}$ equals $\OPT$, it is attained at an endpoint
assignment, and for every $x\in\bar X$
\begin{equation}\label{eq:endpointid}
F(x)-\OPT=\sum_{t}\sum_{v\in E(B_t)}r_t(v)\,\pi_t(v;x)
+\frac12\sum_{i=1}^n(-H_{ii})(x_i-l_i)(u_i-x_i).
\end{equation}
All tables and residuals are computed in
$O\bigl(p\,2^p(N+|\mathcal A|)\bigr)$ table operations on rationals of
bit length polynomial in $I$.
\end{lemma}

\begin{proof}
Since $B_{t'}^{\uparrow}\subseteq B_t$ for a child $t'$, the bracket is a
function of $v_{B_t}$, and $r_t\ge0$ because $M_t$ is its minimum over the
labels of $B_t\setminus B_t^\uparrow$.

\emph{Telescoping.} Let $v\in E([n])$. Every nonroot table $M_t$ occurs in
$\sum_tr_t(v_{B_t})$ once with a plus sign, in the residual of its parent, and
once with a minus sign. Hence $\sum_tr_t(v_{B_t})=\sum_t\phi_t(v_{B_t})-M_{t_0}
=F(v)-M_{t_0}$.

\emph{Expectation.} Fix $x\in\bar X$ and let $Y$ be the independent endpoint
rounding with $\Pr(Y_i=v_i)=e_i(v_i;x_i)/(u_i-l_i)$. Then $\E Y_i=x_i$,
$\operatorname{Var}Y_i=(x_i-l_i)(u_i-x_i)$, $\E Y_iY_j=x_ix_j$ for $i\ne j$,
and $\E Y_i^2=x_i^2+\operatorname{Var}Y_i$. Therefore
$\E F(Y)=F(x)+\frac12\sum_iH_{ii}(x_i-l_i)(u_i-x_i)$, while
$\E r_t(Y_{B_t})=\sum_vr_t(v)\pi_t(v;x)$. Taking expectations in the
telescoping identity gives \eqref{eq:endpointid} with $M_{t_0}$ in place
of $\OPT$.

\emph{Value.} Every term on the right of \eqref{eq:endpointid} is
nonnegative on $\bar X$, so $F\ge M_{t_0}$ on $\bar X\supseteq X$. Trace back:
choose labels of $B_{t_0}$ attaining $M_{t_0}$; for each child $t'$ of a
processed bag $t$, choose the labels of $B_{t'}\setminus B_{t'}^\uparrow$
attaining $M_{t'}(v_{B_{t'}^\uparrow})$. By running intersection, a variable
belongs to $B_t\setminus B_t^\uparrow$ exactly for the bag $t$ nearest the
root among those containing it, and to $B_{t}^{\uparrow}$ for every other bag
containing it. Hence every variable receives exactly one label, and the
resulting endpoint assignment $v^*$ has all residuals zero. By telescoping
$F(v^*)=M_{t_0}$. Both endpoints of every $X_i$ are feasible, so $v^*\in X$
and $\OPT=M_{t_0}$.

\emph{Checking.} The proof uses only the nonnegativity of the residuals and
one endpoint assignment with zero residuals. A checker therefore needs the
tables $M_t$ and that assignment, not the minimization that produced them; any
tables with these two properties give \eqref{eq:endpointid}.

\emph{Cost.} A bag table has at most $2^p$ entries, and children are combined
by addition. By induction on the tree, every message entry is the sum of the
assigned monomials over the subtree at one endpoint assignment. Its bit length
is therefore bounded by the bit length of a common denominator of all endpoint
monomial values plus $\log_2$ of the number of monomials times their maximum
magnitude, both polynomial in $I$.
\end{proof}

\begin{theorem}[All optimizers of a coordinatewise concave quadratic]\label{thm:endpointset}
Under the hypothesis of Lemma~\ref{lem:endpointid}, a point $x\in X$ is a
global optimizer if and only if
\begin{enumerate}
\item[(a)] $x_i\in\{l_i,u_i\}$ for every $i$ with $H_{ii}<0$, and
\item[(b)] $\prod_{i\in B_t}e_i(v_i;x_i)=0$ for every bag $t$ and every
$v\in E(B_t)$ with $r_t(v)>0$.
\end{enumerate}
These are at most $n+N2^p$ factored polynomial equations of degree at most
$p$, together with the original bounds and integrality. They are computed and
checked in $2^{O(p)}\poly(I)$ bit operations, without a growth assumption and
without enumerating interior integer labels or optimal components.
\end{theorem}

\begin{proof}
By \eqref{eq:endpointid}, $x$ is optimal exactly when every nonnegative term
vanishes. The diagonal terms vanish if and only if (a) holds, and
$\pi_t(v;x)=0$ exactly when one factor $e_i(v_i;x_i)$ is zero.
\end{proof}

Condition (b) is a support condition, not a condition on coordinate
projections. For $F=x+y-2xy=x(1-y)+y(1-x)$ on $[0,1]^2$, both optimal
corners $(0,0)$ and $(1,1)$ use every endpoint of every coordinate, yet the
two positive residuals give $x(1-y)=0$ and $y(1-x)=0$, which exclude every
other point. A sum of $m$ disjoint copies has $2^m$ isolated optimizers and
a description with $2m$ equations.

The equations have a finite combinatorial form. For $x\in\bar X$ define its
\emph{face pattern} $\sigma(x)\in\{\mathsf L,\mathsf U,\ast\}^n$ by
$\sigma_i=\mathsf L$ if $x_i=l_i$, $\sigma_i=\mathsf U$ if $x_i=u_i$, and
$\sigma_i=\ast$ otherwise. For a pattern $\tau$ on $B\subseteq[n]$ let
$\operatorname{cube}(\tau)\subseteq E(B)$ consist of the $v$ with $v_i=l_i$
where $\tau_i=\mathsf L$ and $v_i=u_i$ where $\tau_i=\mathsf U$. Then
$\pi_t(v;x)>0$ if and only if $v\in\operatorname{cube}(\sigma(x)_{B_t})$.
Let $\mathcal P_i=\{\mathsf L,\mathsf U\}$, with $\ast$ added when $H_{ii}=0$
and $X_i$ has a point strictly between $l_i$ and $u_i$. Call
$\sigma\in\prod_i\mathcal P_i$ \emph{admissible} if, for every bag,
$\operatorname{cube}(\sigma_{B_t})\subseteq Z_t:=\{v\in E(B_t):r_t(v)=0\}$.

\begin{corollary}[Optimal faces form a tree-structured constraint problem]\label{cor:facecsp}
Under the hypothesis of Lemma~\ref{lem:endpointid}:
\begin{enumerate}
\item[(i)] $x\in X$ is optimal if and only if $\sigma(x)\in\prod_i\mathcal P_i$
and $\sigma(x)$ is admissible.
\item[(ii)] Replacing any $\ast$ of an admissible pattern by $\mathsf L$ or
$\mathsf U$ keeps it admissible, and
$S=\bigcup_{\sigma\text{ admissible}}\bigl(X\cap\bar F(\sigma)\bigr)$, where
$\bar F(\sigma)=\{x:x_i=l_i\text{ if }\sigma_i=\mathsf L,\ x_i=u_i\text{ if }
\sigma_i=\mathsf U\}$ is a closed face of $\bar X$. Distinct admissible
patterns give disjoint nonempty sets $\{x\in X:\sigma(x)=\sigma\}$.
\item[(iii)] Admissible patterns are the solutions of a constraint problem on
the same tree decomposition, with at most three labels per variable and bag
relations
$\mathcal A_t=\{\tau\in\prod_{i\in B_t}\mathcal P_i:\operatorname{cube}(\tau)
\subseteq Z_t\}$. Consequently, in $2^{O(p)}(N+n)\poly(I)$ bit operations
one can decide whether the optimizer is unique, compute the dimension of $S$,
count the admissible patterns, count $|S|$ when $S$ is finite, compute
$\dist(y,S)^2$ and a nearest optimizer for a rational $y$, and minimize a
linear function over $S$.
\end{enumerate}
\end{corollary}

\begin{proof}
(i) Condition (a) of Theorem~\ref{thm:endpointset} says $\sigma_i(x)\ne\ast$
when $H_{ii}<0$; and $\sigma_i(x)=\ast$ is possible only if $X_i$ has an
interior point. Condition (b) says that every $v$ with $\pi_t(v;x)>0$ has
$r_t(v)=0$, that is, $\operatorname{cube}(\sigma(x)_{B_t})\subseteq Z_t$.

(ii) Replacing $\ast$ by an endpoint shrinks every cube. If $\sigma$ is
admissible and $x\in X\cap\bar F(\sigma)$, then $\sigma(x)$ is obtained from
$\sigma$ by replacing some $\ast$ entries by endpoints, so $x$ is optimal by
(i). Conversely $x\in S$ lies in $\bar F(\sigma(x))$. A pattern $\sigma$
determines the set $\{x\in X:\sigma(x)=\sigma\}$, which is nonempty by the
definition of $\mathcal P_i$.

(iii) Each $\mathcal A_t$ is computed by testing at most $3^p$ patterns
against at most $2^p$ residual entries. A constraint problem whose relations
live on the bags of a tree decomposition is solved by nonserial dynamic
programming over $(T,(B_t))$ in the min-sum, max-sum and sum-product
semirings, with $O(p\,3^p)$ operations per bag \cite{Dechter1999}; each unary
cost is charged once, at the bag nearest the root that contains its variable.
The optimizer is unique if and only if exactly one pattern is admissible and
it has no $\ast$, by the disjointness in (ii). The dimension of $S$ is the
maximum number of $\ast$ entries on continuous coordinates. When $S$ is finite,
$|S|=\sum_\sigma\prod_{i:\sigma_i=\ast}|X_i\cap(l_i,u_i)|$. By (ii),
$\dist(y,S)^2=\min_\sigma\sum_ic_i(\sigma_i)$ over admissible $\sigma$, with
$c_i(\mathsf L)=(y_i-l_i)^2$, $c_i(\mathsf U)=(y_i-u_i)^2$ and
$c_i(\ast)=\dist(y_i,X_i)^2$. A linear objective attains its minimum over each
closed face at an endpoint pattern, so it is minimized by the same
recursion over $\{\mathsf L,\mathsf U\}$ labels. All numbers are rationals of
polynomial bit length.
\end{proof}

\begin{remark}[Two-label integer coordinates and multilinear objectives]\label{rem:endpointext}
If $X_i=\{l_i,l_i+1\}$, replacing $\frac12H_{ii}x_i^2$ by
$\frac12H_{ii}\bigl((2l_i+1)x_i-l_i(l_i+1)\bigr)$ does not change $F$ on $X$.
Hence the sign hypothesis is needed only for coordinates with at least three
feasible values; all binary quadratic programs are covered. The proofs also
apply verbatim to $F=\sum_af_a+\sum_i\psi_i(x_i)$ in which each $f_a$ is
affine in each of its variables separately and each $\psi_i$ is concave on
$[l_i,u_i]$. Then $\E f_a(Y)=f_a(x)$, and the diagonal term of
\eqref{eq:endpointid} becomes $\psi_i(x_i)$ minus its chord at $x_i$. That
difference vanishes at an interior point exactly when $\psi_i$ is affine on
$[l_i,u_i]$, which replaces the condition $H_{ii}<0$ in (a).
\end{remark}

Endpoint optimality for nonpositive diagonals is classical; Del Pia and
Khajavirad use it to separate binary-valued from genuinely continuous
variables \cite{DelPiaKhajavirad2026}. Min-sum messages and nonnegative tree
reparameterizations are also standard \cite{Dechter1999,WainwrightJaakkolaWillsky2005}.
The identity \eqref{eq:endpointid} records what these ingredients give for
the entire optimal set. Corollary~\ref{cor:facecsp} shows that the
decomposition also controls the geometry of $S$, not only the optimal value.

\subsection{The diagonal Lagrangian certificate and unknown growth}\label{sec:diagcert}

From here to the end of the section all coordinates are continuous. A point
$\bar x\in X$ is a \emph{box KKT point} if $\partial_iF(\bar x)\ge0$ when
$\bar x_i=l_i$, $\partial_iF(\bar x)\le0$ when $\bar x_i=u_i$, and
$\partial_iF(\bar x)=0$ otherwise. For such a point put
\[
\lambda_i(\bar x)=\frac{2\,|\partial_iF(\bar x)|}{u_i-l_i},\qquad
\Lambda(\bar x)=\operatorname{diag}(\lambda_1(\bar x),\ldots,\lambda_n(\bar x)).
\]
For $\lambda\in\R^n_{\ge0}$ define the Lagrangian dual function
\[
\psi(\lambda)=\inf_{x\in\R^n}\Bigl[F(x)-\frac12\sum_i\lambda_i(x_i-l_i)(u_i-x_i)\Bigr].
\]

\begin{lemma}[Diagonal certificate, full optimal set and invariance]\label{lem:diagcert}
\begin{enumerate}
\item[(i)] For every box KKT point $\bar x$ and every $x\in\R^n$,
\begin{equation}\label{eq:diagid}
F(x)-F(\bar x)=\frac12(x-\bar x)^{\T}\bigl(H+\Lambda(\bar x)\bigr)(x-\bar x)
+\frac12\sum_i\lambda_i(\bar x)(x_i-l_i)(u_i-x_i).
\end{equation}
\item[(ii)] If $H+\Lambda(\bar x)\succeq0$, then $\bar x\in S$ and
\begin{equation}\label{eq:diagset}
S=\bigl\{x\in X:\ (H+\Lambda(\bar x))(x-\bar x)=0,\ \
\lambda_i(\bar x)(x_i-l_i)(u_i-x_i)=0\ \text{for all }i\bigr\}.
\end{equation}
\item[(iii)] $\psi(\lambda)\le\OPT$ for every $\lambda\ge0$, and the
following are equivalent: (a) $H+\Lambda(\bar x)\succeq0$ for some
$\bar x\in S$; (b) $H+\Lambda(\bar x)\succeq0$ for every $\bar x\in S$;
(c) $\psi(\lambda)=\OPT$ for some $\lambda\ge0$. When they hold, the vector
$\lambda(\bar x)$ is the same for every $\bar x\in S$, and it is the unique
maximizer of $\psi$.
\end{enumerate}
\end{lemma}

\begin{proof}
(i) Both sides are quadratic polynomials in $x$ that vanish at $\bar x$.
Their Hessians are $H$ and $H+\Lambda-\Lambda=H$. The $i$th partial
derivative of the right side at $\bar x$ is
$\frac12\lambda_i(\bar x)(u_i+l_i-2\bar x_i)$. This equals
$|\partial_iF(\bar x)|=\partial_iF(\bar x)$ when $\bar x_i=l_i$, equals
$-|\partial_iF(\bar x)|=\partial_iF(\bar x)$ when $\bar x_i=u_i$, and is
$0=\partial_iF(\bar x)$ at an interior coordinate. Two quadratics with equal
value, gradient and Hessian at one point coincide.

(ii) On $X$ both terms on the right of \eqref{eq:diagid} are nonnegative,
so $\bar x$ is optimal, and $x\in X$ is optimal exactly when both vanish. A
positive semidefinite quadratic form vanishes exactly on the kernel of its
matrix.

(iii) For $x\in X$ the bracket in $\psi$ is at most $F(x)$, so
$\psi(\lambda)\le\OPT$. If (a) holds at $\bar x$, then by (i) the bracket with
$\lambda=\lambda(\bar x)$ equals $F(\bar x)+\frac12(x-\bar x)^{\T}(H+\Lambda)(x-\bar x)$,
whose infimum over $\R^n$ is $F(\bar x)=\OPT$; this gives (c). Suppose (c)
holds for some $\lambda$, and let $\bar x\in S$ be arbitrary. Write
$\mathcal L(x)$ for the bracket. Then
$\OPT=\psi(\lambda)\le\mathcal L(\bar x)\le F(\bar x)=\OPT$. Hence
$\lambda_i(\bar x_i-l_i)(u_i-\bar x_i)=0$ for all $i$, and $\bar x$ minimizes
the quadratic $\mathcal L$ over $\R^n$. So
$\nabla^2\mathcal L=H+\operatorname{diag}(\lambda)\succeq0$ and
$0=\partial_i\mathcal L(\bar x)=\partial_iF(\bar x)+\lambda_i\bigl(\bar x_i-\tfrac{l_i+u_i}2\bigr)$.
At an interior coordinate complementarity gives $\lambda_i=0$. At
$\bar x_i=l_i$ the stationarity equation gives
$\lambda_i=2\partial_iF(\bar x)/(u_i-l_i)=\lambda_i(\bar x)$, and similarly at
$\bar x_i=u_i$. Thus $\lambda=\lambda(\bar x)$ and
$H+\Lambda(\bar x)\succeq0$. Since $\bar x\in S$ was arbitrary, this proves (b)
and the uniqueness and invariance statements. Finally (b) implies (a)
because $S\ne\emptyset$.
\end{proof}

\begin{remark}\label{rem:shor}
Call the instances satisfying (a)--(c) the \emph{diagonal certificate class}
$\mathcal C$. Convex box QPs belong to it, since $\Lambda(\bar x)\succeq0$. It also contains indefinite instances with tilted and
disconnected optimal sets. For example $F=(x-y+t/2)^2+\frac18t(1-t)$ on
$[0,1]^3$ has diagonal $(2,2,\frac14)$, the two optimal segments
$\{(v,v,0)\}$ and $\{(v,v+\frac12,1):0\le v\le\frac12\}$, the multiplier
$\Lambda=\operatorname{diag}(0,0,\frac14)$ at every optimizer, and
$H+\Lambda=2(1,-1,\tfrac12)^{\T}(1,-1,\tfrac12)\succeq0$. The function
$\psi$ is the Lagrangian dual for the constraints
$(x_i-l_i)(u_i-x_i)\ge0$, so condition (c) says that this dual, equivalently
the basic Shor semidefinite relaxation, is exact. Sufficient global
optimality conditions of this box-multiplier form are classical
\cite{JeyakumarRubinovWu2006,LiWuQuan2015}, and stronger relaxations with
complete exactness characterizations are studied in \cite{QiuYildirim2024}.
Lemma~\ref{lem:diagcert} is therefore not a new certificate. What we use is
its invariance: whichever optimizer a search finds, the same rational test
accepts it.
\end{remark}



\begin{theorem}[Unknown-growth discovery of a certified optimal set]\label{thm:diagdiscovery}
Let $X$ be a continuous box. Consider the following procedure. For
$K=1,2,4,\ldots$: let $J_K$ be the least $j\ge0$ with $4^j\tau^2\ge4Kns^2$;
run the proximal iteration with guess $K$ for stages $0,\ldots,J_K$, and let
$y=y^{J_K}$. Run REC$(y)$ (Section~\ref{sec:exact}) and let $\bar x$ be a vertex of the
stationary system it solves, if one exists. If $\bar x$ is a box KKT point and
$H+\Lambda(\bar x)\succeq0$, output $\bar x$, $\OPT=F(\bar x)$, $\Lambda(\bar x)$
and the description \eqref{eq:diagset}; otherwise continue with $2K$.
\begin{enumerate}
\item[(a)] Every output is correct, with no assumption on the instance.
\item[(b)] If $F\in\mathcal C$ and $g>0$ is any set-growth constant, the
procedure stops no later than the first guess $K\ge L/g$, so with $K\le2\kappa$.
\item[(c)] Its running time is $f(p,\kappa)\poly(I)$ bit operations with an
absolute polynomial exponent, and the output has bit length $\poly(I)$.
\end{enumerate}
Neither $S$, nor $g$, nor a bound on $\kappa$ is supplied.
\end{theorem}

\begin{proof}
(a) Acceptance requires exactly the hypotheses of
Lemma~\ref{lem:diagcert}(ii), checked in rational arithmetic; the positive
semidefiniteness test is symmetric Gaussian elimination that rejects a
negative pivot and a zero pivot with a nonzero off-diagonal entry in its row.

(b) If $K\ge L/g$, Lemma~\ref{lem:proximal}(iii) gives
$\dist(y,S)^2\le\frac{99}{256}Kns^24^{-J_K}<\tau^2/4$, with $\tau$ as in REC. By
Lemma~\ref{lem:snap}, REC$(y)$ succeeds and every point it returns is
optimal, hence a box KKT point. By Lemma~\ref{lem:diagcert}(iii), (a) implies
(b), so $H+\Lambda(\bar x)\succeq0$ and the test accepts, whichever optimal
component $\bar x$ belongs to.

(c) At most $1+\log_2(2\kappa)$ guesses run. For each,
$J_K+1=O(\log(Kns^2/\tau^2))=\poly(I)+O(\log K)$, and by
Lemma~\ref{lem:proximal}(i),(iv) a stage costs
$O(p(N+|\mathcal A|)K_\theta^p)$ operations on numbers of bit length
$\poly(I+J_K+\mu K_\theta)$, where $\theta^{-1}\le6\sqrt K$ and $\mu=O(\log K)$.
As in Theorem~\ref{thm:approx}, \eqref{eq:logabsorb} absorbs
$\lceil\log_2(n+2)\rceil^p$ into $2p^p(n+2)$, and every factor depending
only on $p$ and $K\le2\kappa$ enters $f$. REC uses original data, so its
linear system has input-size description, and exact linear programming
returns a vertex of polynomial size in polynomial time. The KKT and semidefiniteness tests are polynomial.
\end{proof}

On $\mathcal C$ the procedure always terminates, because a set-growth
constant always exists \cite[Theorem~3.3]{LuoSturm2000}; $\kappa$ governs only
the running time. Outside $\mathcal C$ it may run forever, so an
implementation needs a work limit, and an unsuccessful guess certifies
nothing about membership. For fixed $p$ the bound is polynomial in $\kappa$
and $I$: the table size is $O\bigl((600\sqrt K)^p(C_0p)^p(n+2)\bigr)$, and the
bit lengths are polynomial in $I$ and $\sqrt K\log K$.

\begin{proposition}[Discovery within $\mathcal C$ is hard without the growth parameter]\label{prop:sshard}
For integers $a_1,\ldots,a_m$ and $T$ with $|T|\le A=\sum_i|a_i|$, let
$\sigma_0=0$ and
\[
F(x,\sigma)=\sum_{i=1}^m(\sigma_i-\sigma_{i-1}-a_ix_i)^2+(\sigma_m-T)^2
+\sum_{i=1}^mx_i(1-x_i)
\]
on $x\in[0,1]^m$, $\sigma\in[-A,A]^m$. The bags $\{\sigma_1,x_1\}$ and
$\{\sigma_{i-1},\sigma_i,x_i\}$ $(2\le i\le m)$ form a path decomposition of
size at most $3$. If some $\hat x\in\{0,1\}^m$ has $\sum_ia_i\hat x_i=T$, then
$F\in\mathcal C$, $\OPT=0$, and $S$ consists exactly of the points
$(\hat x,\hat\sigma)$ with $\hat x$ such a solution and $\hat\sigma$ its
partial sums. Otherwise $\OPT>0$. Consequently, unless $\mathrm P=\mathrm{NP}$,
no polynomial-time algorithm returns an exact optimizer for every instance of
$\mathcal C$ with bag size three, even though $\OPT=0$ is known in advance for
these instances, and the ratio $\kappa$ of these instances is not bounded by
any polynomial in their input length.
\end{proposition}

\begin{proof}
$F\ge0$ on the box, and $F=0$ exactly when all squares and all products
$x_i(1-x_i)$ vanish, that is, for binary $x$ with partial sums $\sigma$ and
$\sigma_m=T$; the partial sums lie in $[-A,A]$. If no solution exists, $F>0$
on the compact box. In the feasible case let $\hat z=(\hat x,\hat\sigma)$ be
optimal. The sum of
squares equals $\frac12(z-\hat z)^{\T}M(z-\hat z)$ with $M\succeq0$,
because each square is of an affine function vanishing at $\hat z$. Also
$\sum_ix_i(1-x_i)=\frac12\sum_i2(x_i-0)(1-x_i)$. So
$F(z)-F(\hat z)=\frac12(z-\hat z)^{\T}M(z-\hat z)+\frac12\sum_i2x_i(1-x_i)$.
At $\hat z$ the gradient is $1-2\hat x_i=\pm1$ in $x_i$ and $0$ in
$\sigma$, so $\hat z$ is a box KKT point with $\lambda_i(\hat z)=2$ on $x$ and
$0$ on $\sigma$, and $H+\Lambda(\hat z)=M\succeq0$. Thus $F\in\mathcal C$.
Subset Sum with binary-encoded integers is NP-complete. An algorithm as in the
statement, stopped at its polynomial time bound and followed by a check of its
output, would decide it. By the remark after Theorem~\ref{thm:diagdiscovery},
the same would follow if $\kappa$ were polynomially bounded on these instances.
\end{proof}

The optimal value of an instance of $\mathcal C$ equals $\max\psi$, the value
of a convex relaxation. Proposition~\ref{prop:sshard} shows that the
difficulty is in locating an optimizer, and that a parameter such as $\kappa$
is needed for an efficient exact search. Corollary~\ref{cor:facecsp} has no
counterpart here. With $T=0$ and integers $a_i$ of both signs, the origin is an
optimizer of an instance of $\mathcal C$, and it is the unique optimizer
exactly when no nonempty subset of the $a_i$ sums to zero, which is an
NP-complete question. So even with one optimizer and the full description
\eqref{eq:diagset} in hand, deciding uniqueness is hard.

